% =====================================================================
%  Abstract  (rendered via the class \abstractpage command).
%  Write the numeric blanks last — see TODO markers.
% =====================================================================
\abstractpage{
Continual document understanding requires models to absorb new document types,
schemas, and domains without forgetting previously acquired capabilities.
Multimodal document encoders such as LayoutLMv3 combine text, layout, and visual
pathways, yet it is unknown \emph{which} of these components drives catastrophic
forgetting: existing continual-learning methods treat the network uniformly or
freeze layers heuristically. This thesis presents a per-layer, per-parameter-group
diagnosis of how forgetting distributes across the text, visual, layout, and
cross-modal fusion components of a tri-modal document encoder under continual
sequence-labelling for key-information extraction. Using a controlled,
apples-to-apples protocol that masks individual input modalities of a single
LayoutLMv3 architecture, and measuring representational drift (linear CKA),
per-group Fisher information, and standard forgetting metrics across the
FUNSD\,$\rightarrow$\,CORD\,$\rightarrow$\,SROIE task sequence, we characterise
where forgetting concentrates.
% TODO (after pilot, Week 4): 1--2 sentences stating the dominant-forgetting
% component, e.g. "We find forgetting concentrates in the cross-modal fusion
% layers, while the text stream remains comparatively stable."

Guided by this measurement, we propose \textbf{DocCL}, a mechanism-targeted
continual-learning method that concentrates its regularisation/adaptation on the
dominant-forgetting component identified by the diagnosis. Across three continual
scenarios (class-incremental, domain-incremental, and mixed) and ten baselines
spanning regularisation, replay, prompt-based, and LoRA-based families
(EWC, LwF, ER, DER++, L2P, DualPrompt, CODA-Prompt, O-LoRA, plus naive and joint
bounds), the targeted method
% TODO (after main grid): "reduces forgetting (BWT) by X points / improves
% average entity-F1 by Y" over component-agnostic continual learning while
% adding negligible overhead.
improves continual performance over component-agnostic baselines on the FCS-CL
benchmark. These results establish the per-component diagnostic protocol as a
practical tool for designing continual-learning methods for multimodal document
encoders.\footnote{Source code: \url{https://github.com/hoangthanh283/doccl}
(\texttt{doccl} branch); experiments tracked on Weights \& Biases
(\texttt{doccl-aaai2027}).}

\vspace{0.5em}
\noindent\textbf{Keywords:} continual learning, catastrophic forgetting, document
understanding, key-information extraction, LayoutLMv3, multimodal representation.
}
